\section{Introduction}
\label{sec:introduction}

Breast cancer is among the most commonly diagnosed cancers worldwide and a major cause of cancer-related mortality in women~\cite{sung2021global}. Mammography and breast ultrasound are central to detection and diagnostic assessment. Machine learning and deep learning have achieved strong performance on medical image analysis tasks, including breast imaging classification and computer-aided diagnosis~\cite{litjens2017survey,balkenende2022application}.

Despite these advances, purely neural models can be difficult to interpret because the final prediction need not expose an explicit reasoning process understandable to medical professionals. In high-stakes clinical settings, explainability and trust are especially important~\cite{rudin2019stop,vandervelden2022xai}. One approach is to introduce an intermediate representation of clinically meaningful visual concepts and combine those concepts with explicit domain knowledge~\cite{koh2020concept}.

Neuro-symbolic AI provides a framework for combining neural perception with symbolic reasoning~\cite{garcez2002neural,garcez2023neurosymbolic}. DeepProbLog is particularly relevant because it integrates neural predicates with probabilistic logic programming, enabling reasoning under uncertainty and differentiable learning through the symbolic program~\cite{manhaeve2018deepproblog,manhaeve2021neural}. In this paper, DeepProbLog combines the uncertain visual concepts extracted from a ResNet18 deep architecture's heads with declarative rules inspired by BI-RADS terminology~\cite{acr2013birads}. To achieve this goal, DeepProbLog represents the ResNet18 as a neural predicate inside a probabilistic logic program. Instead of a black-box binary decision malignant/ non-malignant final decision, every prediction decomposes into concepts and  corresponding rule firing probabilities, respectively. Data augmentation is used to conduct a balanced class experimental protocol along with ROI crops, concept warm-up, and a validation-selected high-sensitivity operating point. A quantitative comparison of the DeepProbLog architecture against a pure ResNet18 baseline is provided.

The remaining of the paper is organized as folows: Section II describes the related work, Section III explains the proposed approach, Section IV presents the experimental setup, followed by a presentation and discussion of the obtained results in section V. Finally, the paper is concluded in section VI.